\section{Distributivité de la multiplications par rappor à l'addition et la soustraction}

\newcommand{\myk}{\textcolor{Bleu}{k}}
\newcommand{\mya}{\textcolor{Rouge}{a}}
\newcommand{\myb}{\textcolor{Vert}{b}}
\begin{propbox}
    Pour tous les nombres $\myk$, $\mya$ et $\myb$, on a :
    $$\myk\times(\mya+\myb)=\answer{\myk\times \mya + \myk \times \myb}\text{\hspace{1.5cm} et \hspace{1.5cm}} \myk\times (\mya-\myb) = \answer{\myk\times \mya - \myk\times \myb}$$
\end{propbox}

\begin{vocab}
\begin{itemize}
    \item Le résultat d'une \hlb[RougeClair]{addition}
        (\hlb[VertClair]{soustraction}) s'appelle une \answer[RougeClair]{\bb{somme}}
        (\answer[VertClair]{\bb{différence}}) et les nombres utilisés s'appellent les
        \answer[GrisClair]{termes}.
    \item Le résultat d'une \hlb[BleuClair]{multiplication} s'appelle un
        \answer[BleuClair]{\bb{produit}} et les nombres utilisés s'appellent les
        \answer[BleuClair]{\bb{facteurs}}.
    \item Le résultat d'une \hlb[JauneClair]{division} s'appelle un
        \answer[JauneClair]{\bb{quotient}} et les nombres utilisés s'appellent le
        \answer[JauneClair]{\bb{dividende}} et le \answer[JauneClair]{\bb{diviseur}}.
        (\answer[JauneClair]{$$\text{dividende} \div \text{diviseur} = \text{quotient}$$})
\end{itemize}
\end{vocab}

\begin{vocab}
    La dernière opération réalisée détermine le nom que l'on donne à une
    expression.
\end{vocab}
